\documentclass[12pt,letterpaper]{extarticle}
\newcommand{\packetlevel}{Grades 2–3}
\newcommand{\packetid}{F03-M-v5}
% Shared layout for the Week 3 packets.  Each packet defines \packetlevel and \packetid first.
\usepackage[T1]{fontenc}
\usepackage[utf8]{inputenc}
\usepackage[default]{sourcesanspro}
\usepackage[letterpaper, left=0.5in, right=0.5in, top=0.72in, bottom=0.68in,
            headheight=18pt, headsep=0.2in, footskip=0.34in]{geometry}
\usepackage{xcolor}
\usepackage{tikz}
\usetikzlibrary{arrows.meta}
\usepackage{fancyhdr}
\usepackage{array}

\definecolor{pbgreen}{HTML}{3FAE4A}
\definecolor{pbblue}{HTML}{2F6FD6}
\definecolor{pbred}{HTML}{D8343C}
\definecolor{pbyellow}{HTML}{F5C518}
\definecolor{pbpurple}{HTML}{7D46A8}
\definecolor{pbpink}{HTML}{F07AB4}
\definecolor{pbteal}{HTML}{1EA59B}
\definecolor{pbgray}{HTML}{8C9099}
\definecolor{slotfill}{gray}{0.975}
\definecolor{framecol}{gray}{0.6}

\pagestyle{fancy}
\fancyhf{}
\fancyhead[L]{\normalsize Week 3 / Shuffle machines / \packetlevel}
\fancyfoot[L]{\small Bellingham Math Circle / Week 3 / \packetid}
\fancyfoot[R]{\small \thepage}
\renewcommand{\headrulewidth}{0pt}
\renewcommand{\footrulewidth}{0pt}

\setlength{\parindent}{0pt}
\setlength{\parskip}{0pt}
\raggedright
\pagenumbering{arabic}

\newcommand{\labelfont}{\small}
\newcommand{\cardfont}{\normalsize\bfseries}

\newcounter{prob}
\newcommand{\problem}[1]{\stepcounter{prob}\par\textbf{Problem \theprob:} #1\par}
% a diagram, centred, with space above it
\newcommand{\fig}[2][10pt]{\par\vspace{#1}{\centering\input{fig/#2}\par}}


\begin{document}

One turn: move every block down its own arrow, then slide the whole bottom row straight up into the top row.
When you draw a machine, give every top slot one arrow going out and every bottom slot one arrow coming in.
A block is home when it is back under its own picture.

\vspace{8pt}
\problem{Put each block under its picture and run each machine. For each block, find the first turn after which it is back home. Then find the first turn after which all the blocks are back home at the same time.}

\fig{m-p1a}
\fig[18pt]{m-p1b}

\newpage
\problem{A block's own slot and the slots it visits on its way home make a loop. Without moving any blocks, colour the top slots so that each loop has its own colour, and predict how many turns the machine takes to bring all the blocks home. Then run it to check. Do this for each of the four machines.}

\fig{m-p2a}
\fig[14pt]{m-p2b}

\newpage
\fig[0pt]{m-p2c}
\fig[24pt]{m-p2d}

\newpage
\problem{For each number of turns from 1 to 6, draw a machine with 5 slots that takes exactly that many turns to bring all the blocks home.}

\fig[16pt]{m-p3}

\newpage
\problem{Put each block under its picture on a machine with arrows and do one turn. Move the row of blocks to the empty machine below it without changing its order. Draw arrows on the empty machine so that one turn of it puts every block home. Do this for both machines with arrows.}

\fig{m-p4a}
\fig[24pt]{m-p4a-blank}

\newpage
\fig[0pt]{m-p4b}
\fig[24pt]{m-p4b-blank}

\newpage
\problem{The machines you drew in Problem 4 undo the machines above them. A machine undoes itself if two turns of it put every block home. Draw four different machines with 5 slots that undo themselves. How can you tell from its loops whether a machine undoes itself? Explain.}

\fig[16pt]{m-p5}

\newpage
\problem{Put each block under its picture on machine A and do one turn. Move the row to machine B without changing its order, and do one turn there. Draw the blocks where they end up in the row marked A then B. Then start again with one turn of B followed by one turn of A, and draw the row marked B then A. Do this for both pairs of machines. For which pair are the two rows the same? Explain why.}

\fig[14pt]{m-p6a}

\newpage
\fig[0pt]{m-p6b}

\newpage
\problem{Which numbers of turns can a machine with 6 slots take to bring all the blocks home? Draw a machine for each number you find, and explain why no other number is possible.}

\fig[16pt]{m-p7}

\newpage
\problem{Draw a machine with 7 slots that takes more than 10 turns to bring all the blocks home. What is the most turns a machine with 7 slots can take? Explain.}

\fig[16pt]{m-p8}

\end{document}
